% !TEX root = ../algorithms.tex

\section{Совершенное хеширование. Двухуровневое совершенное хеширование}

\textbf{Задача.} Дано фиксированное конечное множество $S \subseteq U$, $|S| = n$. Требуется построить структуру данных, обеспечивающую поиск любого элемента $x \in U$ за время $O(1)$ в \emph{худшем случае} при использовании $O(n)$ памяти. Структура статическая: операции вставки и удаления не поддерживаются.

\subsection{Однородное совершенное хеширование}

\begin{Lemma}{}{}
	Пусть $\mathcal{H}$ --- универсальное семейство хеш-функций $U \to \{0,\ldots,m-1\}$ и $|S|=n$, $m = n^2$. Тогда при случайном равновероятном выборе $h \in \mathcal{H}$
	\[
	\Pr\bigl(\exists\, x \neq y \in S : h(x) = h(y)\bigr) \;\leq\; \tfrac{1}{2}.
	\]
\end{Lemma}

\begin{proof}
	 Всего неупорядоченных пар $\{x,y\} \subseteq S$, $x \neq y$, ровно $\binom{n}{2}$, и по определению универсального семейства $\Pr(h(x)=h(y)) \leq \tfrac{1}{m} = \tfrac{1}{n^2}$. По неравенству объединения
	\[
	\Pr\bigl(\exists\, x \neq y \in S : h(x) = h(y)\bigr)
	\;\leq\; \binom{n}{2} \cdot \frac{1}{m}
	= \frac{n(n-1)}{2} \cdot \frac{1}{n^2}
	\;\leq\; \frac{1}{2}.
	\]
\end{proof}

\begin{Remark}{}{}
	Это означает, что при многократном независимом выборе $h \in \mathcal{H}$ за ожидаемое число попыток $\leq 2$ найдётся функция $h$, не имеющая коллизий на $S$. Однако размер таблицы $m = n^2$ может быть неприемлемо большим.
\end{Remark}

\subsection{Двухуровневая схема}

Для достижения линейного объёма памяти используется двухуровневая конструкция.

\textbf{Первый уровень.} Выбираем $h\colon U \to \{0,\ldots,n-1\}$ случайно и равновероятно из универсального семейства (то есть на первом уровне $m = n$). Для каждого $k = 0,\ldots,n-1$ полагаем
\[
S_k = \{x \in S : h(x) = k\}, \qquad n_k = |S_k|, \qquad \sum_{k=0}^{n-1} n_k = n.
\]

\textbf{Второй уровень.} В ячейке $T[k]$ первого уровня храним хеш-таблицу второго уровня размера $n_k^2$ с хеш-функцией $h_k$, выбранной случайно и равновероятно из универсального семейства. По предыдущей лемме при каждой повторной попытке выбора $h_k$ с вероятностью $\geq \tfrac{1}{2}$ в таблице $k$-го уровня не будет коллизий.

\begin{figure}[H]
	\centering
	\begin{tikzpicture}[>=Stealth, font=\small]
		\draw (0,0) rectangle (0.6,4);
		\foreach \y in {0.5,1,1.5,2,2.5,3,3.5}
		\draw (0,\y) -- (0.6,\y);
		\node[above] at (0.3,4) {$T$};
		
		\draw[->] (0.6,3.25) -- (1.6,3.25);
		\draw (1.6,3) rectangle (3.6,3.5);
		\node at (2.6,3.25) {$h_2,\; n_2^2$};
		
		\draw[->] (0.6,2.25) -- (1.6,2.25);
		\draw (1.6,2) rectangle (3.0,2.5);
		\node at (2.3,2.25) {$h_5,\; n_5^2$};
		
		\draw[->] (0.6,1.25) -- (1.6,1.25);
		\draw (1.6,1) rectangle (3.3,1.5);
		\node at (2.45,1.25) {$h_8,\; n_8^2$};
	\end{tikzpicture}
	\caption{Двухуровневая схема. Занятая ячейка $T[k]$ первого уровня (размер $n$) ссылается на таблицу второго уровня размера $n_k^2$ с собственной хеш-функцией $h_k$, не имеющей коллизий на $S_k$. Суммарный размер таблиц второго уровня составляет $O(n)$, поэтому вся структура занимает $O(n)$ памяти.}
\end{figure}

\textbf{Поиск.} По запросу $x$: вычислить $k = h(x)$, затем найти $x$ в таблице $S_k$ по адресу $h_k(x)$ --- два вычисления хеш-функции, $O(1)$.

\begin{Proposition}{}{}
	В описанной схеме
	\[
	\mathbf{E}\!\left[\sum_{k=0}^{n-1} n_k^2\right] \;\leq\; 2n.
	\]
\end{Proposition}

\begin{proof}
	Из тождества $n_k^2 = n_k + 2\binom{n_k}{2}$ и равенства $\sum_{k=0}^{n-1} n_k = n$ получаем $\sum_{k=0}^{n-1} n_k^2 = n + 2\sum_{k=0}^{n-1}\binom{n_k}{2}$. Величина $\sum_{k}\binom{n_k}{2}$ равна числу неупорядоченных пар $\{x,y\} \subseteq S$, $x \neq y$, с $h(x) = h(y)$ (числу коллизий первого уровня). По линейности математического ожидания и свойству универсального семейства ($m = n$):
	\[
	\mathbf{E}\!\left[\sum_{k=0}^{n-1}\binom{n_k}{2}\right]
	= \binom{n}{2} \cdot \Pr(h(x)=h(y))
	\;\leq\; \binom{n}{2} \cdot \frac{1}{n}
	= \frac{n-1}{2}.
	\]
	Следовательно, $\mathbf{E}\!\left[\sum_{k=0}^{n-1} n_k^2\right] \leq n + 2 \cdot \tfrac{n-1}{2} = 2n - 1 \leq 2n$.
\end{proof}

\begin{Corollary}{}{}
	По неравенству Маркова, применённому к неотрицательной величине $\sum_k n_k^2$, имеем $\Pr\!\bigl(\sum_k n_k^2 \geq 2\,\mathbf{E}[\sum_k n_k^2]\bigr) \leq \tfrac{1}{2}$. Так как $\mathbf{E}[\sum_k n_k^2] \leq 2n$, при случайном выборе $h$ с вероятностью $\geq \tfrac{1}{2}$ суммарный размер таблиц второго уровня не превышает $4n = O(n)$. Если это не выполнено, выбираем новую функцию $h$ (в среднем $O(1)$ попыток).
\end{Corollary}

\begin{Corollary}{}{}
	Существует хеш-таблица для фиксированного множества $S$, $|S|=n$, занимающая $O(n)$ памяти и обеспечивающая доступ к любому элементу за $O(1)$ в худшем случае.
	
	\emph{Построение:}
	\begin{enumerate}
		\item Выбирать $h \in \mathcal{H}$ (в $U \to \{0,\ldots,n-1\}$) случайно и равновероятно до тех пор, пока $\sum_k n_k^2 \leq 4n$ (ожидаемое число попыток $\leq 2$).
		\item Для каждого $k$ выделить таблицу второго уровня размера $n_k^2$ и выбирать $h_k \in \mathcal{H}$ (в $U \to \{0,\ldots,n_k^2-1\}$) случайно и равновероятно до тех пор, пока коллизий в $S_k$ нет (ожидаемое число попыток $\leq 2$).
		\item Поиск $x$: вычислить $k = h(x)$, затем проверить ячейку $h_k(x)$ таблицы $S_k$ --- одно обращение к памяти, $O(1)$.
	\end{enumerate}
\end{Corollary}